% !PRINTED 65
\weberpage{65}
%<<CONT

ergiebt, so dass das Zeichen \(\Pi(n)\) hier dieselbe Bedeutung hat, wie im ersten Abschnitt (\S~7).

Irgend eine Anordnung des Systems \(\mathfrak A\) bezeichnen wir mit \(\mathfrak A'\), oder ausführlicher, wenn \(\alpha_1,\alpha_2,\ldots,\alpha_n\) die Ziffern \(1,2,\ldots,n\) in irgend einer Reihenfolge bedeuten, mit
\begin{equation}
\mathfrak A'=\alpha_1,\alpha_2,\ldots,\alpha_n.
\tag{3}
\end{equation}

Man kann auf sehr verschiedene Arten aus einer Anordnung eine beliebige andere ableiten, d.~h. eine Permutation ausführen. Unter den verschiedenen Möglichkeiten sind für uns die durch sogenannte \gesperrt{Transpositionen}, d.~h. durch successive Vertauschung von nur zwei Elementen ausgeführten, von besonderem Interesse. Durch mehrere, nach einander ausgeführte Transpositionen lässt sich aus jeder Anordnung, z.~B. aus \(\mathfrak A\), jede andere \(\mathfrak A'\) herleiten. Man kann zu diesem Zweck etwa so verfahren, dass man in \(\mathfrak A\) zunächst das Element \(1\) mit dem, was in \(\mathfrak A'\) an erster Stelle steht, also mit \(\alpha_1\), vertauscht (falls nicht \(\alpha_1=1\) ist), dann, wenn \(\alpha_2\) nicht schon \(=2\) ist, \(2\) mit \(\alpha_2\) u.~s.~f.

Um z.~B. von \((1,2,3,4)\) zu \((4,3,2,1)\) zu gelangen, bildet man die Anordnungen
\[
(1,2,3,4),\qquad (4,2,3,1),\qquad (4,3,2,1).
\]
Bezeichnen wir eine Transposition kurz durch die beiden vertauschten Ziffern, also die Vertauschung von \(1\) mit \(2\) durch \((1,2)\), so haben wir hier nach einander die Transpositionen
\[
(1,4),\quad (2,3)
\]
ausgeführt.

Es ist zu bemerken, dass der Uebergang von einer Anordnung zu einer bestimmten anderen auf unendlich viele verschiedene Arten durch auf einander folgende Transpositionen erreicht werden kann. So geht die Anordnung \((1,2,3,4)\) auch durch die Transpositionen
\[
(1,2),\quad (1,3),\quad (2,4),\quad (1,2)
\]
in \((4,3,2,1)\) über.

\begin{center}
{\large \S.\ 19.}

\vspace{0.5em}
{\large\gesperrt{Permutationen erster und zweiter Art.}}
\end{center}

Die \(\Pi(n)\) Anordnungen von \(n\) Elementen lassen sich nach folgendem Gesichtspunkte in zwei Arten zerlegen.

Aus den \(n\) Elementen unseres Systems lassen sich
\[
{n(n-1)\over2}
\]
und nicht mehr Paare bilden. Wir wollen nun den \(n\) Elementen
%CONT>>
